\documentclass[11pt]{article}
\usepackage[utf8]{inputenc}
\usepackage{amsmath,amssymb}
\usepackage{graphicx}
\usepackage{booktabs}
\usepackage{hyperref}
\usepackage[margin=1in]{geometry}
\usepackage{natbib}
\bibliographystyle{unsrtnat}

\title{Disseminating Cells in Human Oral Tumours Acquire an EMT Cancer Stem Cell State That Is Predictive of Metastasis}
\author{}
\date{}

\begin{document}
\maketitle

\begin{abstract}
The epithelial-to-mesenchymal transition (EMT) is widely proposed to drive metastatic dissemination, yet direct observation of disseminating EMT cancer cells in human tumour specimens has remained elusive because these cells downregulate conventional epithelial markers and become indistinguishable from surrounding stroma. Here we exploit the unique retention of EpCAM on the surface of EMT cancer stem cells (CSCs) to develop an automated 4-colour immunofluorescent imaging and analysis protocol that identifies single disseminating EpCAM$^{+}$Vimentin$^{+}$CD24$^{+}$ cells in the stromal compartment of human oral squamous cell carcinoma specimens. Across 84 tumour specimens and more than 12,000 imaging fields, we show that these triple-positive cells are significantly enriched in the stroma of metastatic tumours compared to non-metastatic tumours, while EpCAM$^{+}$CD24$^{-}$Vimentin$^{+}$ cells show no such correlation, highlighting the specific importance of the CD24$^{+}$ plasticity marker. An artificial neural network trained on the three-marker staining profile predicts metastatic status with 87--89\% cross-validated accuracy, validated on an independent tumour batch. These findings provide the first direct observation of EMT CSCs disseminating from human tumours in vivo and establish a predictive link between single disseminating cells and metastatic outcome.
\end{abstract}

\section{Introduction}

Metastasis is the primary cause of cancer-related mortality, and the epithelial-to-mesenchymal transition has long been proposed as a critical enabling mechanism \citep{thiery2002epithelial,nieto2016emt}. During EMT, epithelial tumour cells acquire mesenchymal traits---including enhanced motility, invasiveness, and resistance to apoptosis---that are thought to facilitate detachment from the primary tumour and dissemination to distant sites \citep{yang2004twist,mani2008mesenchymal}. However, the EMT model of metastasis has been built almost entirely from mouse models and cell line studies, and direct evidence for EMT-driven dissemination in human tumour specimens has been notably absent.

A fundamental obstacle has prevented direct observation of disseminating EMT cells in human tissues: as cancer cells undergo EMT, they downregulate epithelial markers such as pan-keratins and E-cadherin, rendering them indistinguishable from the surrounding mesenchymal stroma \citep{nieto2016emt}. Previous attempts to detect EMT in human tumours have therefore been limited to cells at the tumour margin that retain partial epithelial character, representing only the earliest stages of the transition \citep{brabletz2001variable}. Cells that have completed EMT and entered the stroma become effectively invisible to standard immunohistochemical approaches.

Recent work has identified a subset of EMT cancer stem cells that retain expression of EpCAM and CD24 despite acquiring mesenchymal characteristics \citep{bierie2017integrin,pastushenko2018identification}. These EpCAM$^{+}$CD24$^{+}$ EMT CSCs possess enhanced plasticity, retaining the capacity to undergo mesenchymal-to-epithelial transition (MET) at metastatic sites---a property considered essential for metastatic colonization. Critically, EpCAM is retained on the cell surface during EMT while pan-keratins are lost, making EpCAM a uniquely suitable lineage tracer for cancer cells that have entered the stromal compartment.

Here we combine EpCAM, Vimentin (a mesenchymal marker upregulated during EMT), and CD24 (a marker of EMT CSC plasticity) in an automated 4-colour immunofluorescent imaging protocol to identify and quantify single disseminating EMT CSCs in the stroma of 84 human oral squamous cell carcinoma (OSCC) specimens. We demonstrate that these cells are significantly enriched in metastatic tumours and develop a machine learning classifier that predicts metastatic status from the three-marker profile.

\section{Results}

\subsection{EpCAM$^{+}$Vimentin$^{+}$CD24$^{+}$ staining identifies EMT cancer stem cells}

We first validated the triple-positive staining profile using the CA1 oral squamous cell carcinoma cell line and its EMT-stem derivative sub-line, which has been previously characterized as possessing EMT CSC properties with retained plasticity. Immunofluorescent staining for EpCAM (yellow), Vimentin (red), and CD24 (green), with DAPI nuclear counterstain (blue), revealed that EpCAM$^{+}$Vimentin$^{+}$CD24$^{+}$ triple-positive cells constituted 41.0\% of the EMT-stem sub-line compared to only 2.1\% of the parental CA1 line (Figure~1A--B,E; $p < 0.05$, two-tailed Student's $t$-test). Normal keratinocytes and cancer-associated fibroblasts showed no triple-positive cells, confirming specificity (Figure~1, Supplement~1).

A parallel staining panel replacing EpCAM with pan-keratin showed no enrichment of pan-keratin$^{+}$Vimentin$^{+}$CD24$^{+}$ cells in the EMT-stem sub-line (Figure~1C--D), confirming that pan-keratins are lost during EMT while EpCAM is retained. This result establishes EpCAM as the critical epithelial marker for detecting cancer cells that have undergone EMT and entered the mesenchymal compartment.

\subsection{Automated imaging identifies disseminating cells in human tumour stroma}

We applied the validated 4-colour staining protocol to 84 human OSCC specimens, stratified by metastatic status, using automated whole-slide imaging that captured more than 12,000 imaging fields (Figure~2A--D). Image segmentation was performed to identify individual cells and quantify marker co-expression at single-cell resolution. An EpCAM density map was generated for each specimen to delineate the tumour body (high EpCAM density) from the surrounding stroma (low EpCAM density), enabling specific quantification of cells that had disseminated beyond the tumour boundary (Figure~2E).

Single EpCAM$^{+}$ cells detected in the stromal compartment represent tumour cells that have left the main tumour mass. Among these disseminating cells, we quantified the proportion expressing the full EpCAM$^{+}$Vimentin$^{+}$CD24$^{+}$ triple-positive profile indicative of EMT CSC status.

\subsection{Triple-positive disseminating cells are enriched in metastatic tumours}

Comparison of metastatic and non-metastatic specimens revealed a significant enrichment of single EpCAM$^{+}$Vimentin$^{+}$CD24$^{+}$ cells in the stroma of metastatic tumours ($p < 0.05$; Figure~2). This enrichment was specific to the triple-positive population: EpCAM$^{+}$Vimentin$^{+}$CD24$^{-}$ cells in the stroma showed no significant correlation with metastatic status, and the pan-keratin$^{+}$Vimentin$^{+}$CD24$^{+}$ control staining showed no enrichment.

This finding demonstrates that the CD24$^{+}$ plasticity marker is specifically required, alongside EpCAM and Vimentin, to identify the subset of disseminating EMT cells that is associated with metastatic disease. Cells that have undergone EMT but lack CD24 expression---and therefore potentially lack the capacity for MET at metastatic sites---do not correlate with metastatic outcome.

\subsection{Machine learning predicts metastasis from staining profiles}

To test whether the spatial distribution of the three markers contains sufficient information to predict metastatic status, we trained an artificial neural network (ANN) on grey-level intensities of EpCAM, Vimentin, and CD24 in imaging tiles from tumour batch~1 (Figure~3A).

Using 10-fold cross-validation, the ANN achieved a classification accuracy of 87--89\% in distinguishing metastatic from non-metastatic tumour specimens based on the EpCAM/Vimentin/CD24 staining panel (Figure~3B). The same architecture trained on the pan-keratin/Vimentin/CD24 control panel achieved substantially lower accuracy (Figure~3C), confirming that the predictive information resides specifically in the EpCAM-based staining profile.

The model trained on batch~1 was then applied to batch~2 as an independent test set, maintaining high classification accuracy (Figure~3). This cross-batch generalization demonstrates that the predictive signal is robust and not driven by batch-specific technical variation.

\section{Discussion}

Our study provides the first direct observation of EMT cancer stem cells disseminating from human tumours in vivo and establishes a quantitative link between these cells and metastatic outcome. By exploiting the unique retention of EpCAM on the surface of EMT CSCs, we overcame the fundamental barrier that has prevented detection of disseminating EMT cells in human tissue sections---their loss of conventional epithelial markers.

The specific requirement for CD24 positivity is noteworthy. CD24 marks a subset of EMT CSCs that retain the capacity for mesenchymal-to-epithelial transition, a property considered essential for metastatic colonization \citep{bierie2017integrin}. Our finding that only CD24$^{+}$ disseminating cells---not CD24$^{-}$ cells that have also undergone EMT---correlate with metastasis supports the hypothesis that plasticity, rather than EMT per se, is the critical determinant of metastatic competence \citep{pastushenko2018identification,jolly2019implications}.

The 87--89\% cross-validated accuracy of the neural network classifier suggests clinical potential for this approach. The three-marker immunofluorescent staining protocol uses standard reagents applicable to routine formalin-fixed paraffin-embedded specimens, and the automated imaging and analysis pipeline could be integrated into pathological workflows. However, prospective validation in larger, multi-institutional cohorts will be essential before clinical implementation.

Several limitations merit consideration. The current study examined oral squamous cell carcinoma; generalization to other tumour types requires validation. The binary classification (metastatic vs. non-metastatic) does not capture the temporal dynamics of dissemination, and longitudinal studies tracking patient outcomes from the time of primary tumour resection would strengthen the prognostic value of the approach. Additionally, while our imaging identifies single disseminating cells, it cannot determine whether these cells will successfully complete the metastatic cascade.

\section{Methods}

\subsection{Cell line validation}
The CA1 OSCC cell line and its EMT-stem derivative sub-line were stained with two panels: (1) EpCAM/Vimentin/CD24/DAPI and (2) pan-keratin/Vimentin/CD24/DAPI. Normal keratinocytes and cancer-associated fibroblasts served as negative controls. Quantification was performed using two-tailed Student's $t$-tests with 95\% confidence intervals.

\subsection{Human tumour specimens}
Eighty-four formalin-fixed paraffin-embedded OSCC specimens, stratified by metastatic status, were processed in two independent batches. Four-colour immunofluorescent staining (EpCAM, Vimentin, CD24, DAPI) was applied, and entire histopathological slides were captured using automated imaging, yielding more than 12,000 imaging fields.

\subsection{Image analysis}
Image segmentation identified individual cells and quantified marker co-expression. EpCAM density maps were generated to delineate the tumour body from the surrounding stroma. Disseminating cells were defined as single EpCAM$^{+}$ cells located in the stromal compartment beyond the tumour body boundary.

\subsection{Machine learning}
A supervised artificial neural network was trained on grey-level intensities of the three markers (EpCAM, Vimentin, CD24) in imaging tiles. Training was performed on batch~1 with 10-fold cross-validation. The trained model was independently validated on batch~2.

\bibliography{references}

\end{document}
